%\vspace{-2pt}
\section{Introduction}
\label{sec:intro}
\vspace{-5pt}
Reconstructing three-dimensional scenes from sparse image inputs remains a significant challenge in computer vision, with applications spanning robotics, autonomous systems, and augmented/virtual reality. Traditional deep multi-view stereo (MVS) based methods, such as MVSNet~\cite{yao2018mvsnet}, have established a foundation for this task by estimating depth maps through the construction of 3D cost volumes within the camera's viewing frustum. Subsequent refinements, including approaches by Yang et al.~\cite{yang2020cost}, Gu et al.~\cite{gu2020cascade}, Wang et al.~\cite{wang2021patchmatchnet}, and Ding et al.~\cite{ding2022transmvsnet}, have improved depth accuracy. However, these methods often require extensive post-processing, such as depth map filtering, and struggle with low-texture regions, noise sensitivity, and incomplete data, particularly with sparse image views. Additionally, they can not enable stand-alone novel view synthesis (NVS).

Neural implicit representation techniques have emerged as powerful alternatives, providing high-fidelity 3D reconstructions by implicitly representing surfaces through neural Signed Distance Functions (SDF)~\cite{park2019deepsdf}. Advances in volume rendering~\cite{yariv2021volume, wang2021neus, oechsle2021unisurf, niemeyer2020differentiable, darmon2022improving} have enabled smoother and more detailed reconstructions by directly optimizing scene geometry and radiance from multi-view images. Recently, 3D Gaussian Splatting (3DGS)~\cite{kerbl20233d} introduced the use of Gaussian primitives for faster real-time, photorealistic NVS. While this method provides significant improvements, obtaining continuous 3D surfaces from these primitives remains challenging. Solutions have emerged to improve surface extraction. For instance,  SuGaR~\cite{guedon2024sugar} enhances surface alignment via a post-hock processing. Furthermore, 2D Gaussian Splatting (2DGS)~\cite{huang20242d}  improves reconstruction accuracy via a tuned primitive representation and an improved rendering algorithm.

Despite these advances, most current methods, especially the scene specific test-time optimization ones, still face limitations including high computational demands,  extensive input view requirements, and limited generalization across different scenes. To address these issues, generalizable 3D reconstruction and NVS models~\cite{johari2022geonerf, xu2023wavenerf, chen2021mvsnerf, ren2023volrecon, long2022sparseneus, liang2024retr, jena2024geotransfer, na2024uforecon} aim to utilize image features to predict radiance and signed distance fields using feed-forward networks trained on large datasets. However, being volumetric rendering based methods, they inherit notably slow inference speeds. Recent work has introduced feed-forward optimization-free Gaussian Splatting models~\cite{liu2024fast, chen2024mvsplat, charatan2024pixelsplat, szymanowicz2024splatter, zheng2024gps} that use pixel-aligned 3D Gaussian primitives~\cite{kerbl20233d} for fast novel view synthesis from sparse views, by virtue of the tile-based rasterization and fast GPU sorting algorithms used in the splatting process. These methods offer cross-scene generalization even with limited image data. However, they are not fit for feedforward surface reconstruction because of their NVS purposed output 3D representation. In this respect, we find (in Sec. \ref{sec:abl}) that SOTA generalizable GS method \cite{liu2024fast} for novel view synthesis fails to give coherent reconstructions on TSDF fusion of the rendered depth maps.  

% Ins

% In this paper, we build on these approaches to address sparse 3D reconstruction challenges by employing the large-scale pretrained foundation models, DinoV2~\cite{oquab2023dinov2} and MASt3R~\cite{leroy2024grounding}. 

In this paper, we tackle the problem of fast multi-view shape reconstruction from sparse input views. In this regard, we build a generalizable feed-forward model that regresses 2D Gaussian splatting parameters (instead of 3D primitives) which allows us to reconstruct notably faster than the previous generalizable 3D reconstruction methods  relying mostly on implicit volumetric rendering~\cite{johari2022geonerf, xu2023wavenerf, chen2021mvsnerf, ren2023volrecon, long2022sparseneus, liang2024retr, jena2024geotransfer, na2024uforecon}. Inspired by MVSFormer++~\cite{cao2024mvsformer++}, which leverages $2$D foundation model DINOv2~\cite{oquab2023dinov2} for feature encoding with cross-view attention and enhances cost volume regularization for state-of-the-art multi-view Stereo, we extend the MVSGaussian~\cite{liu2024fast} framework to utilize monocular or multi-view features extracted from $2$D and $3$D foundation models for the task of $3$D reconstruction and novel view synthesis. Our method investigates the use of rich 2D semantic features from DINOv2~\cite{oquab2023dinov2} and pairwise feature maps from MASt3R~\cite{leroy2024grounding} encoding dense pairwise correspondences between input images to predict $2$DGS parameters, and demonstrates state-of-the-art results in both $3$D reconstruction and novel view synthesis on sparse-view setups. In summary, our contributions are as follows:
\vspace{-1pt}
\begin{itemize}
\item We propose the first generalizable, feed-forward approach for sparse-view novel view synthesis and 3D reconstruction using 2D Gaussian splatting. %by leveraging large-scale pretrained foundation models, DinoV2~\cite{oquab2023dinov2} and MASt3R~\cite{leroy2024grounding}.
% \item We extend the MVSGaussian~\cite{liu2024fast} framework to integrate both monocular and multi-view features obtained from foundation models. 
\item Inspired by MVSFormer++\cite{cao2024mvsformer++}, which leverages DinoV2\cite{oquab2023dinov2} for feature encoding and cost volume regularization, we investigate the impact of incorporating 2D semantic monocular features from DinoV2~\cite{oquab2023dinov2} and dense pairwise correspondence features from MASt3R~\cite{leroy2024grounding} for predicting 2D Gaussian Splatting parameters.
\item We conduct extensive experiments and ablations on the DTU dataset~\cite{aanaes2016large}, demonstrating that our approach, powered by MASt3R~\cite{leroy2024grounding} features achieves state-of-the-art results in both 3D reconstruction and novel view synthesis, and fast inference.
\end{itemize}